  %%%%%%%%%%%%%%%%%%%%%%%%%%%%%%%%%%%%%%%%%%%%%%%%%%%%%%%%%%%%
% 0.8 METHODS OF PROOF
% The Composition of Advanced Mathematics
%%%%%%%%%%%%%%%%%%%%%%%%%%%%%%%%%%%%%%%%%%%%%%%%%%%%%%%%%%%%

\section{Methods of Proof}
\label{sec:methods-of-proof}

\prerequisite{
Familiarity with propositions, implications, quantifiers, and logical
equivalence from \cref{sec:elementary-logic}, together with sets, relations,
and functions from the preceding sections.
}

\learningobjectives{
By the end of this section, the reader should be able to:
\begin{itemize}
    \item distinguish mathematical evidence from mathematical proof;
    \item identify hypotheses and conclusions in theorem statements;
    \item use definitions as tools in proofs;
    \item construct direct proofs;
    \item prove statements by contraposition and contradiction;
    \item organize proofs by cases;
    \item prove existence and uniqueness statements;
    \item distinguish constructive and nonconstructive existence proofs;
    \item disprove universal statements using counterexamples;
    \item prove biconditional statements;
    \item prove set containments and equalities;
    \item prove basic statements involving functions;
    \item use mathematical induction and strong induction;
    \item recognize recursive definitions;
    \item select an appropriate proof strategy for a given statement;
    \item identify common logical and structural errors in proofs.
\end{itemize}
}

%%%%%%%%%%%%%%%%%%%%%%%%%%%%%%%%%%%%%%%%%%%%%%%%%%%%%%%%%%%%
\subsection{What Is a Mathematical Proof?}
\label{subsec:what-is-proof}
%%%%%%%%%%%%%%%%%%%%%%%%%%%%%%%%%%%%%%%%%%%%%%%%%%%%%%%%%%%%

A mathematical proof is a logically valid argument establishing that a
mathematical statement must be true under the stated assumptions.

A proof does more than provide evidence.

It explains why the conclusion follows necessarily from definitions,
assumptions, previously established results, and valid logical reasoning.

\begin{definitionbox}{Mathematical Proof}
A \emph{mathematical proof} is a finite sequence of logically justified
statements that establishes the truth of a mathematical claim.
\end{definitionbox}

%%%%%%%%%%%%%%%%%%%%%%%%%%%%%%%%%%%%%%%%%%%%%%%%%%%%%%%%%%%%
\subsection{Evidence versus Proof}
\label{subsec:evidence-versus-proof}
%%%%%%%%%%%%%%%%%%%%%%%%%%%%%%%%%%%%%%%%%%%%%%%%%%%%%%%%%%%%

Examples can provide evidence for a mathematical statement, but examples alone
cannot prove a universal claim.

Suppose one suspects that

\[ n^2+n \]

is even for every integer \(n\).

Checking

\[ 1^2+1=2,\qquad 2^2+2=6,\qquad 3^2+3=12 \]

provides evidence.

However, infinitely many integers remain unchecked.

A proof must establish the result for an arbitrary integer.

%%%%%%%%%%%%%%%%%%%%%%%%%%%%%%%%%%%%%%%%%%%%%%%%%%%%%%%%%%%%
\subsection{The Role of Definitions}
\label{subsec:definitions-in-proof}
%%%%%%%%%%%%%%%%%%%%%%%%%%%%%%%%%%%%%%%%%%%%%%%%%%%%%%%%%%%%

Definitions are among the most important tools in proof writing.

For example, the statement

\begin{center}
``\(n\) is even''
\end{center}

has the precise meaning

\[ n=2k \]

for some integer \(k\).

Similarly,

\begin{center}
``\(n\) is odd''
\end{center}

means

\[ n=2k+1 \]

for some integer \(k\).

A proof often begins by replacing terminology with its formal definition.

%%%%%%%%%%%%%%%%%%%%%%%%%%%%%%%%%%%%%%%%%%%%%%%%%%%%%%%%%%%%
\subsection{Theorem Structure}
\label{subsec:theorem-structure}
%%%%%%%%%%%%%%%%%%%%%%%%%%%%%%%%%%%%%%%%%%%%%%%%%%%%%%%%%%%%

Many theorems have the logical form

\[ P\to Q. \]

Here \(P\) is the hypothesis and \(Q\) is the conclusion.

For example:

\begin{theorem}
If \(n\) is an even integer, then \(n^2\) is even.
\end{theorem}

The hypothesis is

\begin{center}
\(n\) is even.
\end{center}

The conclusion is

\begin{center}
\(n^2\) is even.
\end{center}

%%%%%%%%%%%%%%%%%%%%%%%%%%%%%%%%%%%%%%%%%%%%%%%%%%%%%%%%%%%%
\subsection{Beginning a Proof}
\label{subsec:beginning-proof}
%%%%%%%%%%%%%%%%%%%%%%%%%%%%%%%%%%%%%%%%%%%%%%%%%%%%%%%%%%%%

When proving a universal implication such as

\[ \forall x\in A,\;P(x)\to Q(x), \]

a standard beginning is:

\begin{enumerate}
    \item let \(x\in A\) be arbitrary;
    \item assume \(P(x)\);
    \item use the hypothesis and known facts;
    \item establish \(Q(x)\).
\end{enumerate}

The word \emph{arbitrary} matters.

It indicates that the argument does not depend on a special choice of \(x\).

%%%%%%%%%%%%%%%%%%%%%%%%%%%%%%%%%%%%%%%%%%%%%%%%%%%%%%%%%%%%
\subsection{Direct Proof}
\label{subsec:direct-proof}
%%%%%%%%%%%%%%%%%%%%%%%%%%%%%%%%%%%%%%%%%%%%%%%%%%%%%%%%%%%%

\begin{definitionbox}{Direct Proof}
A \emph{direct proof} of
\[ P\to Q \]
begins by assuming \(P\) and proceeds logically until \(Q\) is established.
\end{definitionbox}

This is often the most natural proof method when the hypothesis can be
translated directly into useful algebraic or structural information.

%%%%%%%%%%%%%%%%%%%%%%%%%%%%%%%%%%%%%%%%%%%%%%%%%%%%%%%%%%%%
\subsection{Worked Example: Direct Proof}
\label{subsec:worked-direct-proof}
%%%%%%%%%%%%%%%%%%%%%%%%%%%%%%%%%%%%%%%%%%%%%%%%%%%%%%%%%%%%

\begin{workedexamplebox}{The Sum of Two Even Integers Is Even}

Prove that if \(a\) and \(b\) are even integers, then \(a+b\) is even.

\begin{tabularx}{\textwidth}{@{}>{\raggedright\arraybackslash}p{0.42\textwidth}X@{}}
Use the definition of even. & \(a=2m,\qquad b=2n\) for some \(m,n\in\Z\)\\
Add the two integers. & \(a+b=2m+2n\)\\
Factor out \(2\). & \(a+b=2(m+n)\)\\
Use closure of integers under addition. & \(m+n\in\Z\)
\end{tabularx}

Since \(a+b\) is \(2\) times an integer,

\begin{answerbox}
\(a+b\) is even.
\end{answerbox}
\end{workedexamplebox}

%%%%%%%%%%%%%%%%%%%%%%%%%%%%%%%%%%%%%%%%%%%%%%%%%%%%%%%%%%%%
\subsection{Writing the Arbitrary Element}
\label{subsec:arbitrary-element}
%%%%%%%%%%%%%%%%%%%%%%%%%%%%%%%%%%%%%%%%%%%%%%%%%%%%%%%%%%%%

A proof involving ``every'' or ``for all'' usually requires an arbitrary
element.

For example, to prove

\[ A\subseteq B, \]

one begins:

\begin{center}
Let \(x\in A\) be arbitrary.
\end{center}

The goal is then to show

\[ x\in B. \]

Because \(x\) was arbitrary, the result holds for every element of \(A\).

%%%%%%%%%%%%%%%%%%%%%%%%%%%%%%%%%%%%%%%%%%%%%%%%%%%%%%%%%%%%
\subsection{Proof by Cases}
\label{subsec:proof-by-cases}
%%%%%%%%%%%%%%%%%%%%%%%%%%%%%%%%%%%%%%%%%%%%%%%%%%%%%%%%%%%%

Sometimes the hypothesis naturally divides into several possibilities.

\begin{definitionbox}{Proof by Cases}
A \emph{proof by cases} divides the problem into exhaustive cases and proves
the desired conclusion separately in each case.
\end{definitionbox}

The cases must collectively cover every possibility permitted by the
hypothesis.

%%%%%%%%%%%%%%%%%%%%%%%%%%%%%%%%%%%%%%%%%%%%%%%%%%%%%%%%%%%%
\subsection{Worked Example: Proof by Cases}
\label{subsec:worked-proof-cases}
%%%%%%%%%%%%%%%%%%%%%%%%%%%%%%%%%%%%%%%%%%%%%%%%%%%%%%%%%%%%

\begin{workedexamplebox}{The Square of an Integer Has Predictable Parity}

Prove that for every integer \(n\), \(n^2+n\) is even.

Every integer is either even or odd.

\textbf{Case 1: \(n\) is even.}

Then

\[ n=2k \]

for some \(k\in\Z\). Hence

\[ n^2+n=(2k)^2+2k=4k^2+2k=2(2k^2+k), \]

which is even.

\textbf{Case 2: \(n\) is odd.}

Then

\[ n=2k+1. \]

Thus

\[ n^2+n=(2k+1)^2+(2k+1)=4k^2+6k+2=2(2k^2+3k+1), \]

which is also even.

\begin{answerbox}
Therefore \(n^2+n\) is even for every integer \(n\).
\end{answerbox}
\end{workedexamplebox}

%%%%%%%%%%%%%%%%%%%%%%%%%%%%%%%%%%%%%%%%%%%%%%%%%%%%%%%%%%%%
\subsection{Exhaustive Cases}
\label{subsec:exhaustive-cases}
%%%%%%%%%%%%%%%%%%%%%%%%%%%%%%%%%%%%%%%%%%%%%%%%%%%%%%%%%%%%

A case proof is valid only when all possibilities are covered.

For integers, the cases

\begin{center}
even or odd
\end{center}

are exhaustive.

Similarly, a real number satisfies exactly one of

\[ x<0,\qquad x=0,\qquad x>0. \]

Leaving out one possibility creates an incomplete proof.

%%%%%%%%%%%%%%%%%%%%%%%%%%%%%%%%%%%%%%%%%%%%%%%%%%%%%%%%%%%%
\subsection{Proof by Contraposition}
\label{subsec:proof-contraposition}
%%%%%%%%%%%%%%%%%%%%%%%%%%%%%%%%%%%%%%%%%%%%%%%%%%%%%%%%%%%%

Recall the logical equivalence

\[ P\to Q\equiv\neg Q\to\neg P. \]

\begin{definitionbox}{Proof by Contraposition}
To prove
\[ P\to Q, \]
one may instead prove the logically equivalent contrapositive
\[ \neg Q\to\neg P. \]
\end{definitionbox}

This strategy is useful when the negation of the conclusion gives more usable
information than the original hypothesis.

%%%%%%%%%%%%%%%%%%%%%%%%%%%%%%%%%%%%%%%%%%%%%%%%%%%%%%%%%%%%
\subsection{Worked Example: Contraposition}
\label{subsec:worked-contraposition}
%%%%%%%%%%%%%%%%%%%%%%%%%%%%%%%%%%%%%%%%%%%%%%%%%%%%%%%%%%%%

\begin{workedexamplebox}{If \(n^2\) Is Odd, Then \(n\) Is Odd}

Instead of proving the statement directly, prove its contrapositive:

\begin{center}
If \(n\) is even, then \(n^2\) is even.
\end{center}

\begin{tabularx}{\textwidth}{@{}>{\raggedright\arraybackslash}p{0.42\textwidth}X@{}}
Assume \(n\) is even. & \(n=2k\) for some \(k\in\Z\)\\
Square both sides. & \(n^2=(2k)^2=4k^2\)\\
Factor out \(2\). & \(n^2=2(2k^2)\)\\
Since \(2k^2\in\Z\). & \(n^2\) is even
\end{tabularx}

The contrapositive is true, so the original implication is true.

\begin{answerbox}
If \(n^2\) is odd, then \(n\) is odd.
\end{answerbox}
\end{workedexamplebox}

%%%%%%%%%%%%%%%%%%%%%%%%%%%%%%%%%%%%%%%%%%%%%%%%%%%%%%%%%%%%
\subsection{Choosing Contraposition}
\label{subsec:choosing-contraposition}
%%%%%%%%%%%%%%%%%%%%%%%%%%%%%%%%%%%%%%%%%%%%%%%%%%%%%%%%%%%%

Contraposition is especially useful when the conclusion says that something
has a property such as:

\begin{itemize}
    \item being irrational;
    \item being nonzero;
    \item being injective;
    \item not belonging to a set;
    \item failing to satisfy some condition.
\end{itemize}

Negating such conclusions often produces a concrete assumption that is easier
to manipulate.

%%%%%%%%%%%%%%%%%%%%%%%%%%%%%%%%%%%%%%%%%%%%%%%%%%%%%%%%%%%%
\subsection{Proof by Contradiction}
\label{subsec:proof-contradiction}
%%%%%%%%%%%%%%%%%%%%%%%%%%%%%%%%%%%%%%%%%%%%%%%%%%%%%%%%%%%%

The symbol

\[ \bot \]

is called the \emph{contradiction symbol} or \emph{bottom symbol} and is often
read

\begin{center}
``contradiction.''
\end{center}

\begin{definitionbox}{Proof by Contradiction}
To prove a statement \(P\) by contradiction, assume \(\neg P\) and logically
derive an impossibility or contradiction.
\end{definitionbox}

Symbolically,

\[ \neg P\Longrightarrow\bot. \]

Therefore the assumption \(\neg P\) must be false, and \(P\) must be true.

%%%%%%%%%%%%%%%%%%%%%%%%%%%%%%%%%%%%%%%%%%%%%%%%%%%%%%%%%%%%
\subsection{What Counts as a Contradiction?}
\label{subsec:what-is-contradiction-proof}
%%%%%%%%%%%%%%%%%%%%%%%%%%%%%%%%%%%%%%%%%%%%%%%%%%%%%%%%%%%%

A contradiction may appear as:

\[ 0=1, \]

\[ a<a, \]

\[ n\text{ is both even and odd}, \]

or a statement that directly conflicts with an earlier assumption.

The contradiction must arise logically from the negation of the claim and
accepted mathematical facts.

%%%%%%%%%%%%%%%%%%%%%%%%%%%%%%%%%%%%%%%%%%%%%%%%%%%%%%%%%%%%
\subsection{Worked Example: Irrationality of \(\sqrt2\)}
\label{subsec:worked-sqrt2-contradiction}
%%%%%%%%%%%%%%%%%%%%%%%%%%%%%%%%%%%%%%%%%%%%%%%%%%%%%%%%%%%%

\begin{workedexamplebox}{Prove \(\sqrt2\) Is Irrational}

Assume for contradiction that

\[ \sqrt2=\frac ab \]

for integers \(a,b\) with \(b\neq0\), where \(a/b\) is in lowest terms.

\begin{tabularx}{\textwidth}{@{}>{\raggedright\arraybackslash}p{0.42\textwidth}X@{}}
Square both sides. & \(2=\frac{a^2}{b^2}\)\\
Clear the denominator. & \(a^2=2b^2\)\\
Therefore \(a^2\) is even. & Hence \(a\) is even\\
Write \(a=2k\). & \(4k^2=2b^2\)\\
Simplify. & \(b^2=2k^2\)\\
Therefore \(b\) is even. & \(2\mid b\)
\end{tabularx}

Thus both \(a\) and \(b\) are divisible by \(2\), contradicting the assumption
that \(a/b\) was in lowest terms.

\begin{answerbox}
\[ \sqrt2\notin\Q \]
\end{answerbox}
\end{workedexamplebox}

%%%%%%%%%%%%%%%%%%%%%%%%%%%%%%%%%%%%%%%%%%%%%%%%%%%%%%%%%%%%
\subsection{Contraposition versus Contradiction}
\label{subsec:contraposition-vs-contradiction}
%%%%%%%%%%%%%%%%%%%%%%%%%%%%%%%%%%%%%%%%%%%%%%%%%%%%%%%%%%%%

These methods are related but distinct.

To prove

\[ P\to Q \]

by contraposition, prove

\[ \neg Q\to\neg P. \]

To prove the same statement by contradiction, assume

\[ P\land\neg Q \]

and derive a contradiction.

Contraposition often produces a more direct argument.

Contradiction is useful when the negation of the desired statement has strong
structural consequences.

%%%%%%%%%%%%%%%%%%%%%%%%%%%%%%%%%%%%%%%%%%%%%%%%%%%%%%%%%%%%
\subsection{Existence Proofs}
\label{subsec:existence-proofs}
%%%%%%%%%%%%%%%%%%%%%%%%%%%%%%%%%%%%%%%%%%%%%%%%%%%%%%%%%%%%

An existence statement has the form

\[ \exists x\in A,\;P(x). \]

To prove it, one must establish that at least one suitable object exists.

There are two broad approaches:

\begin{itemize}
    \item constructive existence;
    \item nonconstructive existence.
\end{itemize}

%%%%%%%%%%%%%%%%%%%%%%%%%%%%%%%%%%%%%%%%%%%%%%%%%%%%%%%%%%%%
\subsection{Constructive Existence}
\label{subsec:constructive-existence}
%%%%%%%%%%%%%%%%%%%%%%%%%%%%%%%%%%%%%%%%%%%%%%%%%%%%%%%%%%%%

\begin{definitionbox}{Constructive Proof}
A \emph{constructive existence proof} establishes existence by explicitly
producing an object with the required property.
\end{definitionbox}

For example, to prove

\begin{center}
There exists an even prime number,
\end{center}

one may simply exhibit

\[ 2. \]

The number \(2\) is prime and even.

%%%%%%%%%%%%%%%%%%%%%%%%%%%%%%%%%%%%%%%%%%%%%%%%%%%%%%%%%%%%
\subsection{Worked Example: Constructive Existence}
\label{subsec:worked-constructive-existence}
%%%%%%%%%%%%%%%%%%%%%%%%%%%%%%%%%%%%%%%%%%%%%%%%%%%%%%%%%%%%

\begin{workedexamplebox}{Existence of an Irrational Sum}

Show that there exist irrational numbers \(a,b\) such that \(a+b\) is rational.

\begin{tabularx}{\textwidth}{@{}>{\raggedright\arraybackslash}p{0.42\textwidth}X@{}}
Choose an irrational number. & \(a=\sqrt2\)\\
Choose its negative. & \(b=-\sqrt2\)\\
Both are irrational. & \(a,b\notin\Q\)\\
Add them. & \(a+b=0\)
\end{tabularx}

\begin{answerbox}
Such irrational numbers exist because \(\sqrt2+(-\sqrt2)=0\in\Q\).
\end{answerbox}
\end{workedexamplebox}

%%%%%%%%%%%%%%%%%%%%%%%%%%%%%%%%%%%%%%%%%%%%%%%%%%%%%%%%%%%%
\subsection{Nonconstructive Existence}
\label{subsec:nonconstructive-existence}
%%%%%%%%%%%%%%%%%%%%%%%%%%%%%%%%%%%%%%%%%%%%%%%%%%%%%%%%%%%%

\begin{definitionbox}{Nonconstructive Proof}
A \emph{nonconstructive existence proof} establishes that an object must
exist without necessarily identifying a specific example.
\end{definitionbox}

Such proofs often use contradiction, counting arguments, completeness
properties, or logical alternatives.

Nonconstructive methods become increasingly common in advanced mathematics.

%%%%%%%%%%%%%%%%%%%%%%%%%%%%%%%%%%%%%%%%%%%%%%%%%%%%%%%%%%%%
\subsection{Uniqueness Proofs}
\label{subsec:uniqueness-proofs}
%%%%%%%%%%%%%%%%%%%%%%%%%%%%%%%%%%%%%%%%%%%%%%%%%%%%%%%%%%%%

A uniqueness statement asserts that exactly one object satisfies a property.

Symbolically,

\[ \exists!x\in A,\;P(x). \]

A complete proof usually contains two parts:

\begin{enumerate}
    \item existence: show at least one such object exists;
    \item uniqueness: show any two such objects must be equal.
\end{enumerate}

%%%%%%%%%%%%%%%%%%%%%%%%%%%%%%%%%%%%%%%%%%%%%%%%%%%%%%%%%%%%
\subsection{Worked Example: Existence and Uniqueness}
\label{subsec:worked-existence-uniqueness}
%%%%%%%%%%%%%%%%%%%%%%%%%%%%%%%%%%%%%%%%%%%%%%%%%%%%%%%%%%%%

\begin{workedexamplebox}{Unique Solution of a Linear Equation}

Prove that for \(a\neq0\), the equation

\[ ax+b=0 \]

has exactly one real solution.

\textbf{Existence.}

Choose

\[ x=-\frac ba. \]

Then

\[ a\left(-\frac ba\right)+b=-b+b=0. \]

\textbf{Uniqueness.}

Suppose \(x_1\) and \(x_2\) both satisfy the equation.

\begin{tabularx}{\textwidth}{@{}>{\raggedright\arraybackslash}p{0.42\textwidth}X@{}}
Write both equations. & \(ax_1+b=0,\qquad ax_2+b=0\)\\
Subtract. & \(a(x_1-x_2)=0\)\\
Use \(a\neq0\). & \(x_1-x_2=0\)\\
Conclude. & \(x_1=x_2\)
\end{tabularx}

\begin{answerbox}
\[ \exists!x\in\R\text{ such that }ax+b=0,\qquad x=-\frac ba. \]
\end{answerbox}
\end{workedexamplebox}

%%%%%%%%%%%%%%%%%%%%%%%%%%%%%%%%%%%%%%%%%%%%%%%%%%%%%%%%%%%%
\subsection{Disproving Statements}
\label{subsec:disproving-statements}
%%%%%%%%%%%%%%%%%%%%%%%%%%%%%%%%%%%%%%%%%%%%%%%%%%%%%%%%%%%%

Not every mathematical claim is true.

A universal statement

\[ \forall x\in A,\;P(x) \]

is disproved by finding one element \(a\in A\) for which \(P(a)\) is false.

Such an object is called a counterexample.

%%%%%%%%%%%%%%%%%%%%%%%%%%%%%%%%%%%%%%%%%%%%%%%%%%%%%%%%%%%%
\subsection{Counterexamples}
\label{subsec:proof-counterexamples}
%%%%%%%%%%%%%%%%%%%%%%%%%%%%%%%%%%%%%%%%%%%%%%%%%%%%%%%%%%%%

\begin{definitionbox}{Counterexample}
A \emph{counterexample} is a specific example showing that a universal
statement is false.
\end{definitionbox}

For example, consider the claim

\begin{center}
Every prime number is odd.
\end{center}

The number

\[ 2 \]

is prime and even.

Therefore \(2\) is a counterexample.

%%%%%%%%%%%%%%%%%%%%%%%%%%%%%%%%%%%%%%%%%%%%%%%%%%%%%%%%%%%%
\subsection{One Counterexample Is Enough}
\label{subsec:one-counterexample}
%%%%%%%%%%%%%%%%%%%%%%%%%%%%%%%%%%%%%%%%%%%%%%%%%%%%%%%%%%%%

To prove a universal claim, one must handle every possible case.

To disprove a universal claim, one counterexample is sufficient.

Symbolically,

\[ \neg(\forall x\,P(x))\equiv\exists x\,\neg P(x). \]

This is the logical foundation of the counterexample method.

%%%%%%%%%%%%%%%%%%%%%%%%%%%%%%%%%%%%%%%%%%%%%%%%%%%%%%%%%%%%
\subsection{Biconditional Proofs}
\label{subsec:biconditional-proofs}
%%%%%%%%%%%%%%%%%%%%%%%%%%%%%%%%%%%%%%%%%%%%%%%%%%%%%%%%%%%%

A statement of the form

\[ P\leftrightarrow Q \]

requires proving both directions:

\[ P\to Q \]

and

\[ Q\to P. \]

These are often labeled:

\begin{center}
Forward direction
\end{center}

and

\begin{center}
Reverse direction.
\end{center}

%%%%%%%%%%%%%%%%%%%%%%%%%%%%%%%%%%%%%%%%%%%%%%%%%%%%%%%%%%%%
\subsection{Worked Example: Biconditional}
\label{subsec:worked-biconditional}
%%%%%%%%%%%%%%%%%%%%%%%%%%%%%%%%%%%%%%%%%%%%%%%%%%%%%%%%%%%%

\begin{workedexamplebox}{An Integer Is Even iff Its Square Is Even}

Prove that an integer \(n\) is even if and only if \(n^2\) is even.

\textbf{Forward direction.}

Assume \(n\) is even. Then

\[ n=2k \]

for some \(k\in\Z\), so

\[ n^2=4k^2=2(2k^2). \]

Therefore \(n^2\) is even.

\textbf{Reverse direction.}

Assume \(n^2\) is even.

The contrapositive of the desired implication is:

\begin{center}
If \(n\) is odd, then \(n^2\) is odd.
\end{center}

Write

\[ n=2k+1. \]

Then

\[ n^2=4k^2+4k+1=2(2k^2+2k)+1, \]

which is odd.

Therefore \(n^2\) even implies \(n\) even.

\begin{answerbox}
\[ n\text{ is even}\Longleftrightarrow n^2\text{ is even}. \]
\end{answerbox}
\end{workedexamplebox}

%%%%%%%%%%%%%%%%%%%%%%%%%%%%%%%%%%%%%%%%%%%%%%%%%%%%%%%%%%%%
\subsection{Proofs of Set Containment}
\label{subsec:proof-set-containment}
%%%%%%%%%%%%%%%%%%%%%%%%%%%%%%%%%%%%%%%%%%%%%%%%%%%%%%%%%%%%

To prove

\[ A\subseteq B, \]

take an arbitrary element

\[ x\in A \]

and show

\[ x\in B. \]

This translates the subset statement into an implication about membership.

%%%%%%%%%%%%%%%%%%%%%%%%%%%%%%%%%%%%%%%%%%%%%%%%%%%%%%%%%%%%
\subsection{Proofs of Set Equality}
\label{subsec:proof-set-equality}
%%%%%%%%%%%%%%%%%%%%%%%%%%%%%%%%%%%%%%%%%%%%%%%%%%%%%%%%%%%%

To prove

\[ A=B, \]

a standard method is double containment:

\[ A\subseteq B \]

and

\[ B\subseteq A. \]

This approach was introduced in \cref{sec:elementary-set-theory} and is one of
the most common proof patterns in higher mathematics.

%%%%%%%%%%%%%%%%%%%%%%%%%%%%%%%%%%%%%%%%%%%%%%%%%%%%%%%%%%%%
\subsection{Worked Example: Set Equality}
\label{subsec:worked-set-equality-proof}
%%%%%%%%%%%%%%%%%%%%%%%%%%%%%%%%%%%%%%%%%%%%%%%%%%%%%%%%%%%%

\begin{workedexamplebox}{Prove a De Morgan Law}

Prove

\[ (A\cup B)^c=A^c\cap B^c. \]

Let \(x\) be arbitrary.

\begin{tabularx}{\textwidth}{@{}>{\raggedright\arraybackslash}p{0.42\textwidth}X@{}}
Start on the left. & \(x\in(A\cup B)^c\)\\
Use complement. & \(x\notin A\cup B\)\\
Use the definition of union. & \(x\notin A\) and \(x\notin B\)\\
Use complement again. & \(x\in A^c\) and \(x\in B^c\)\\
Use intersection. & \(x\in A^c\cap B^c\)
\end{tabularx}

The argument is reversible, so membership in either set is equivalent.

\begin{answerbox}
\[ (A\cup B)^c=A^c\cap B^c \]
\end{answerbox}
\end{workedexamplebox}

%%%%%%%%%%%%%%%%%%%%%%%%%%%%%%%%%%%%%%%%%%%%%%%%%%%%%%%%%%%%
\subsection{Proving Function Injectivity}
\label{subsec:proof-injectivity}
%%%%%%%%%%%%%%%%%%%%%%%%%%%%%%%%%%%%%%%%%%%%%%%%%%%%%%%%%%%%

To prove that

\[ f:A\to B \]

is injective, a common strategy is:

\begin{enumerate}
    \item assume \(f(x_1)=f(x_2)\);
    \item manipulate the equality;
    \item prove \(x_1=x_2\).
\end{enumerate}

This follows directly from the definition of injectivity.

%%%%%%%%%%%%%%%%%%%%%%%%%%%%%%%%%%%%%%%%%%%%%%%%%%%%%%%%%%%%
\subsection{Worked Example: Injectivity}
\label{subsec:worked-proof-injectivity}
%%%%%%%%%%%%%%%%%%%%%%%%%%%%%%%%%%%%%%%%%%%%%%%%%%%%%%%%%%%%

\begin{workedexamplebox}{Prove a Linear Function Is Injective}

Let

\[ f:\R\to\R,\qquad f(x)=5x-4. \]

\begin{tabularx}{\textwidth}{@{}>{\raggedright\arraybackslash}p{0.42\textwidth}X@{}}
Assume equal outputs. & \(f(x_1)=f(x_2)\)\\
Substitute the formula. & \(5x_1-4=5x_2-4\)\\
Add \(4\). & \(5x_1=5x_2\)\\
Divide by \(5\). & \(x_1=x_2\)
\end{tabularx}

\begin{answerbox}
\(f\) is injective.
\end{answerbox}
\end{workedexamplebox}

%%%%%%%%%%%%%%%%%%%%%%%%%%%%%%%%%%%%%%%%%%%%%%%%%%%%%%%%%%%%
\subsection{Proving Function Surjectivity}
\label{subsec:proof-surjectivity}
%%%%%%%%%%%%%%%%%%%%%%%%%%%%%%%%%%%%%%%%%%%%%%%%%%%%%%%%%%%%

To prove that

\[ f:A\to B \]

is surjective:

\begin{enumerate}
    \item let \(y\in B\) be arbitrary;
    \item find or construct an \(x\in A\);
    \item show \(f(x)=y\).
\end{enumerate}

This matches the definition

\[ \forall y\in B,\;\exists x\in A\text{ such that }f(x)=y. \]

%%%%%%%%%%%%%%%%%%%%%%%%%%%%%%%%%%%%%%%%%%%%%%%%%%%%%%%%%%%%
\subsection{Worked Example: Surjectivity}
\label{subsec:worked-proof-surjectivity}
%%%%%%%%%%%%%%%%%%%%%%%%%%%%%%%%%%%%%%%%%%%%%%%%%%%%%%%%%%%%

\begin{workedexamplebox}{Prove a Linear Function Is Surjective}

Let

\[ f:\R\to\R,\qquad f(x)=5x-4. \]

Take arbitrary

\[ y\in\R. \]

We seek \(x\) satisfying

\[ 5x-4=y. \]

\begin{tabularx}{\textwidth}{@{}>{\raggedright\arraybackslash}p{0.42\textwidth}X@{}}
Add \(4\). & \(5x=y+4\)\\
Divide by \(5\). & \(x=\frac{y+4}{5}\)\\
This value is real. & \(x\in\R\)\\
Substitute. & \(f(x)=y\)
\end{tabularx}

\begin{answerbox}
\(f\) is surjective.
\end{answerbox}
\end{workedexamplebox}

%%%%%%%%%%%%%%%%%%%%%%%%%%%%%%%%%%%%%%%%%%%%%%%%%%%%%%%%%%%%
\subsection{Mathematical Induction}
\label{subsec:mathematical-induction}
%%%%%%%%%%%%%%%%%%%%%%%%%%%%%%%%%%%%%%%%%%%%%%%%%%%%%%%%%%%%

Many statements concern every natural number.

Mathematical induction provides a systematic method for proving such
statements.

\begin{definitionbox}{Mathematical Induction}
Suppose \(P(n)\) is a statement defined for natural numbers. To prove
\[ P(n)\text{ for all }n\geq n_0, \]
it is sufficient to show:
\begin{enumerate}
    \item \(P(n_0)\) is true;
    \item for every \(k\geq n_0\), \(P(k)\to P(k+1)\).
\end{enumerate}
\end{definitionbox}

The first step is called the \emph{base case}.

The second is called the \emph{inductive step}.

%%%%%%%%%%%%%%%%%%%%%%%%%%%%%%%%%%%%%%%%%%%%%%%%%%%%%%%%%%%%
\subsection{The Inductive Hypothesis}
\label{subsec:inductive-hypothesis}
%%%%%%%%%%%%%%%%%%%%%%%%%%%%%%%%%%%%%%%%%%%%%%%%%%%%%%%%%%%%

In the inductive step, one assumes

\[ P(k) \]

for an arbitrary integer \(k\geq n_0\).

This assumption is called the \emph{inductive hypothesis}.

The goal is then to prove

\[ P(k+1). \]

One does not assume that \(P(k+1)\) is true.

That is precisely what must be established.

%%%%%%%%%%%%%%%%%%%%%%%%%%%%%%%%%%%%%%%%%%%%%%%%%%%%%%%%%%%%
\subsection{Why Induction Works}
\label{subsec:why-induction-works}
%%%%%%%%%%%%%%%%%%%%%%%%%%%%%%%%%%%%%%%%%%%%%%%%%%%%%%%%%%%%

Induction can be visualized as a chain.

The base case establishes the first link.

The inductive step proves that whenever one link holds, the next one follows:

\[ P(n_0)\to P(n_0+1)\to P(n_0+2)\to\cdots. \]

Thus every subsequent statement is reached.

%%%%%%%%%%%%%%%%%%%%%%%%%%%%%%%%%%%%%%%%%%%%%%%%%%%%%%%%%%%%
\subsection{Worked Example: Sum of the First \(n\) Integers}
\label{subsec:worked-induction-sum}
%%%%%%%%%%%%%%%%%%%%%%%%%%%%%%%%%%%%%%%%%%%%%%%%%%%%%%%%%%%%

\begin{workedexamplebox}{Prove a Summation Formula}

Prove that for every \(n\in\N\),

\[ 1+2+\cdots+n=\frac{n(n+1)}2. \]

\textbf{Base case: \(n=1\).}

\[ 1=\frac{1(1+1)}2=1. \]

So the statement holds for \(n=1\).

\textbf{Inductive hypothesis.}

Assume for some \(k\in\N\),

\[ 1+2+\cdots+k=\frac{k(k+1)}2. \]

\textbf{Inductive step.}

Then

\[
\begin{aligned}
1+2+\cdots+k+(k+1)
&=\frac{k(k+1)}2+(k+1)\\
&=(k+1)\left(\frac{k}{2}+1\right)\\
&=\frac{(k+1)(k+2)}2.
\end{aligned}
\]

This is exactly the desired formula with \(n=k+1\).

\begin{answerbox}
\[ 1+2+\cdots+n=\frac{n(n+1)}2 \qquad\text{for all }n\in\N. \]
\end{answerbox}
\end{workedexamplebox}

%%%%%%%%%%%%%%%%%%%%%%%%%%%%%%%%%%%%%%%%%%%%%%%%%%%%%%%%%%%%
\subsection{Summation Notation}
\label{subsec:summation-notation-proof}
%%%%%%%%%%%%%%%%%%%%%%%%%%%%%%%%%%%%%%%%%%%%%%%%%%%%%%%%%%%%

The symbol

\[ \sum \]

is the \emph{Greek capital sigma}, pronounced ``SIG-muh.''

In mathematics, it denotes a sum.

For example,

\[ \sum_{k=1}^{n}k \]

is read

\begin{center}
``the sum from \(k=1\) to \(n\) of \(k\).''
\end{center}

Thus,

\[ \sum_{k=1}^{n}k=1+2+\cdots+n. \]

The previous induction result may therefore be written

\[ \sum_{k=1}^{n}k=\frac{n(n+1)}2. \]

%%%%%%%%%%%%%%%%%%%%%%%%%%%%%%%%%%%%%%%%%%%%%%%%%%%%%%%%%%%%
\subsection{Induction with a Different Starting Value}
\label{subsec:induction-starting-index}
%%%%%%%%%%%%%%%%%%%%%%%%%%%%%%%%%%%%%%%%%%%%%%%%%%%%%%%%%%%%

Induction does not have to begin at \(1\).

If a statement is intended for every integer

\[ n\geq5, \]

then the natural base case is

\[ n=5. \]

The inductive step then assumes the result for \(k\geq5\) and proves it for
\(k+1\).

%%%%%%%%%%%%%%%%%%%%%%%%%%%%%%%%%%%%%%%%%%%%%%%%%%%%%%%%%%%%
\subsection{Strong Induction}
\label{subsec:strong-induction}
%%%%%%%%%%%%%%%%%%%%%%%%%%%%%%%%%%%%%%%%%%%%%%%%%%%%%%%%%%%%

\begin{definitionbox}{Strong Induction}
In \emph{strong induction}, the inductive step assumes that
\[ P(n_0),P(n_0+1),\ldots,P(k) \]
are all true and uses these assumptions to prove \(P(k+1)\).
\end{definitionbox}

Strong induction is logically equivalent in power to ordinary induction, but
it is often more convenient when the next case depends on several earlier
cases.

%%%%%%%%%%%%%%%%%%%%%%%%%%%%%%%%%%%%%%%%%%%%%%%%%%%%%%%%%%%%
\subsection{Example Motivation for Strong Induction}
\label{subsec:strong-induction-motivation}
%%%%%%%%%%%%%%%%%%%%%%%%%%%%%%%%%%%%%%%%%%%%%%%%%%%%%%%%%%%%

Suppose a sequence satisfies

\[ a_n=a_{n-1}+a_{n-2}. \]

To reason about \(a_{k+1}\), one may need information about both

\[ a_k \]

and

\[ a_{k-1}. \]

Strong induction naturally allows access to all previous cases.

%%%%%%%%%%%%%%%%%%%%%%%%%%%%%%%%%%%%%%%%%%%%%%%%%%%%%%%%%%%%
\subsection{Worked Example: Prime Factorization Existence}
\label{subsec:worked-strong-induction}
%%%%%%%%%%%%%%%%%%%%%%%%%%%%%%%%%%%%%%%%%%%%%%%%%%%%%%%%%%%%

\begin{workedexamplebox}{Every Integer Greater Than \(1\) Is a Product of Primes}

We prove the statement by strong induction.

\textbf{Base case: \(n=2\).}

The number \(2\) is prime, so it is already a product consisting of one prime.

\textbf{Inductive hypothesis.}

Assume every integer \(m\) satisfying

\[ 2\leq m\leq k \]

is a product of primes.

\textbf{Inductive step.}

Consider \(k+1\).

If \(k+1\) is prime, the result is immediate.

If \(k+1\) is composite, then

\[ k+1=ab \]

for integers

\[ 2\leq a,b\leq k. \]

By the inductive hypothesis, both \(a\) and \(b\) are products of primes.

Therefore their product \(k+1\) is also a product of primes.

\begin{answerbox}
Every integer \(n\geq2\) can be written as a product of primes.
\end{answerbox}
\end{workedexamplebox}

%%%%%%%%%%%%%%%%%%%%%%%%%%%%%%%%%%%%%%%%%%%%%%%%%%%%%%%%%%%%
\subsection{Recursive Definitions}
\label{subsec:recursive-definitions}
%%%%%%%%%%%%%%%%%%%%%%%%%%%%%%%%%%%%%%%%%%%%%%%%%%%%%%%%%%%%

A definition may describe an object using earlier instances of the same
object.

Such a definition is called \emph{recursive}.

For example, define the factorial function by

\[ 0!=1 \]

and

\[ n!=(n)(n-1)! \qquad\text{for }n\geq1. \]

The symbol

\[ ! \]

is called the \emph{factorial symbol} in this context.

The expression

\[ n! \]

is read

\begin{center}
``\(n\) factorial.''
\end{center}

%%%%%%%%%%%%%%%%%%%%%%%%%%%%%%%%%%%%%%%%%%%%%%%%%%%%%%%%%%%%
\subsection{Factorials}
\label{subsec:factorials-proof}
%%%%%%%%%%%%%%%%%%%%%%%%%%%%%%%%%%%%%%%%%%%%%%%%%%%%%%%%%%%%

For positive integers,

\[ n!=n(n-1)(n-2)\cdots2\cdot1. \]

Examples include

\[ 1!=1,\qquad 2!=2,\qquad 3!=6,\qquad 4!=24. \]

Recursive definitions and induction are closely connected because induction
often proves properties of recursively defined objects.

%%%%%%%%%%%%%%%%%%%%%%%%%%%%%%%%%%%%%%%%%%%%%%%%%%%%%%%%%%%%
\subsection{Structural Pattern of a Proof}
\label{subsec:structural-pattern-proof}
%%%%%%%%%%%%%%%%%%%%%%%%%%%%%%%%%%%%%%%%%%%%%%%%%%%%%%%%%%%%

A well-written proof generally contains:

\begin{enumerate}
    \item a clear statement of the assumptions;
    \item appropriate definitions;
    \item a logically connected sequence of steps;
    \item justification for important transitions;
    \item a clear conclusion.
\end{enumerate}

A proof should communicate why the result follows, not merely display symbolic
manipulations.

%%%%%%%%%%%%%%%%%%%%%%%%%%%%%%%%%%%%%%%%%%%%%%%%%%%%%%%%%%%%
\subsection{Choosing a Proof Method}
\label{subsec:choosing-proof-method}
%%%%%%%%%%%%%%%%%%%%%%%%%%%%%%%%%%%%%%%%%%%%%%%%%%%%%%%%%%%%

The logical form of the statement often suggests a method.

\[
\begin{array}{c|l}
\text{Statement Form} & \text{Common Strategy}\\ \hline
P\to Q & \text{direct proof}\\
P\to Q & \text{contraposition}\\
P & \text{contradiction}\\
P\leftrightarrow Q & \text{prove both directions}\\
\forall x\,P(x) & \text{arbitrary element}\\
\exists x\,P(x) & \text{construct an example}\\
\exists!x\,P(x) & \text{existence + uniqueness}\\
A\subseteq B & \text{element argument}\\
A=B & \text{double containment}\\
\forall n\in\N\,P(n) & \text{induction may apply}
\end{array}
\]

These are guidelines rather than rigid rules.

%%%%%%%%%%%%%%%%%%%%%%%%%%%%%%%%%%%%%%%%%%%%%%%%%%%%%%%%%%%%
\subsection{Work Backward, Write Forward}
\label{subsec:work-backward-write-forward}
%%%%%%%%%%%%%%%%%%%%%%%%%%%%%%%%%%%%%%%%%%%%%%%%%%%%%%%%%%%%

When discovering a proof, it is often useful to work backward from the desired
conclusion.

Suppose the goal is

\[ a^2=b^2. \]

One may ask what would be sufficient to guarantee this.

Perhaps

\[ a=b \]

would be enough.

This can help discover a proof.

However, the final written proof should generally proceed logically from the
hypotheses toward the conclusion.

%%%%%%%%%%%%%%%%%%%%%%%%%%%%%%%%%%%%%%%%%%%%%%%%%%%%%%%%%%%%
\subsection{Common Proof Errors}
\label{subsec:common-proof-errors}
%%%%%%%%%%%%%%%%%%%%%%%%%%%%%%%%%%%%%%%%%%%%%%%%%%%%%%%%%%%%

\subsubsection{Proving by Example}

Checking

\[ n=1,2,3,4 \]

does not prove a statement for every integer.

Examples can suggest a theorem but do not establish a universal claim.

%%%%%%%%%%%%%%%%%%%%%%%%%%%%%%%%%%%%%%%%%%%%%%%%%%%%%%%%%%%%
\subsubsection{Assuming What Must Be Proved}
\label{subsec:circular-proof}
%%%%%%%%%%%%%%%%%%%%%%%%%%%%%%%%%%%%%%%%%%%%%%%%%%%%%%%%%%%%

A proof cannot assume its desired conclusion unless that assumption is part of
a legitimate contradiction argument.

Using the conclusion as justification for itself is called \emph{circular
reasoning}.

%%%%%%%%%%%%%%%%%%%%%%%%%%%%%%%%%%%%%%%%%%%%%%%%%%%%%%%%%%%%
\subsubsection{Using the Converse}
\label{subsec:proof-error-converse}
%%%%%%%%%%%%%%%%%%%%%%%%%%%%%%%%%%%%%%%%%%%%%%%%%%%%%%%%%%%%

From

\[ P\to Q \]

one cannot automatically use

\[ Q\to P. \]

The converse must be established independently.

%%%%%%%%%%%%%%%%%%%%%%%%%%%%%%%%%%%%%%%%%%%%%%%%%%%%%%%%%%%%
\subsubsection{Dividing by a Possibly Zero Quantity}
\label{subsec:proof-error-zero-division}
%%%%%%%%%%%%%%%%%%%%%%%%%%%%%%%%%%%%%%%%%%%%%%%%%%%%%%%%%%%%

A step such as

\[ ax=ay\Longrightarrow x=y \]

requires

\[ a\neq0. \]

Without that condition, the conclusion may fail.

Every cancellation or division step must respect the relevant restrictions.

%%%%%%%%%%%%%%%%%%%%%%%%%%%%%%%%%%%%%%%%%%%%%%%%%%%%%%%%%%%%
\subsubsection{Using an Unproved Special Case as a General Result}
\label{subsec:proof-error-generalization}
%%%%%%%%%%%%%%%%%%%%%%%%%%%%%%%%%%%%%%%%%%%%%%%%%%%%%%%%%%%%

A statement shown for a particular object cannot automatically be generalized.

The proof must explain why the argument applies to an arbitrary object in the
required domain.

%%%%%%%%%%%%%%%%%%%%%%%%%%%%%%%%%%%%%%%%%%%%%%%%%%%%%%%%%%%%
\subsubsection{Incomplete Case Analysis}
\label{subsec:proof-error-cases}
%%%%%%%%%%%%%%%%%%%%%%%%%%%%%%%%%%%%%%%%%%%%%%%%%%%%%%%%%%%%

A proof by cases must cover every possibility.

If a real-number proof considers only

\[ x>0 \]

and

\[ x<0, \]

it has omitted

\[ x=0. \]

%%%%%%%%%%%%%%%%%%%%%%%%%%%%%%%%%%%%%%%%%%%%%%%%%%%%%%%%%%%%
\subsubsection{Misusing the Inductive Hypothesis}
\label{subsec:proof-error-induction}
%%%%%%%%%%%%%%%%%%%%%%%%%%%%%%%%%%%%%%%%%%%%%%%%%%%%%%%%%%%%

In induction, one may assume \(P(k)\) while proving \(P(k+1)\).

One may not assume

\[ P(k+1) \]

itself.

Doing so would make the argument circular.

%%%%%%%%%%%%%%%%%%%%%%%%%%%%%%%%%%%%%%%%%%%%%%%%%%%%%%%%%%%%
\subsection{Worked Example: Finding the Flaw}
\label{subsec:worked-false-proof}
%%%%%%%%%%%%%%%%%%%%%%%%%%%%%%%%%%%%%%%%%%%%%%%%%%%%%%%%%%%%

\begin{workedexamplebox}{A False Cancellation Argument}

Consider the following argument.

Suppose

\[ a=b. \]

Then

\[ a^2=ab. \]

Subtract \(b^2\):

\[ a^2-b^2=ab-b^2. \]

Factor:

\[ (a-b)(a+b)=b(a-b). \]

Cancel \(a-b\):

\[ a+b=b. \]

Since \(a=b\),

\[ 2b=b, \]

so

\[ 2=1. \]

Where is the error?

\begin{tabularx}{\textwidth}{@{}>{\raggedright\arraybackslash}p{0.42\textwidth}X@{}}
The assumption is \(a=b\). & \(a-b=0\)\\
The proof cancels \(a-b\). & This divides both sides by \(0\)\\
Division by zero is undefined. & The cancellation is invalid
\end{tabularx}

\begin{answerbox}
The argument fails because it divides by \(a-b=0\).
\end{answerbox}
\end{workedexamplebox}

%%%%%%%%%%%%%%%%%%%%%%%%%%%%%%%%%%%%%%%%%%%%%%%%%%%%%%%%%%%%
\subsection{Proof Writing and Communication}
\label{subsec:proof-writing-communication}
%%%%%%%%%%%%%%%%%%%%%%%%%%%%%%%%%%%%%%%%%%%%%%%%%%%%%%%%%%%%

A proof is written for a reader.

Good proof writing should make the logical structure visible.

Useful phrases include:

\begin{itemize}
    \item ``Let \(x\) be arbitrary.''
    \item ``Assume that...''
    \item ``By definition...''
    \item ``Since...''
    \item ``Therefore...''
    \item ``Conversely...''
    \item ``Suppose for contradiction that...''
    \item ``By the inductive hypothesis...''
\end{itemize}

Symbols should support the reasoning rather than replace complete explanations.

%%%%%%%%%%%%%%%%%%%%%%%%%%%%%%%%%%%%%%%%%%%%%%%%%%%%%%%%%%%%
\subsection{Exercises}
\label{subsec:proof-exercises}
%%%%%%%%%%%%%%%%%%%%%%%%%%%%%%%%%%%%%%%%%%%%%%%%%%%%%%%%%%%%

\begin{exercise}[1]
Prove directly that the sum of two odd integers is even.
\end{exercise}

\begin{exercise}[1]
Prove directly that the product of two even integers is even.
\end{exercise}

\begin{exercise}[1]
Prove that the product of an even integer and any integer is even.
\end{exercise}

\begin{exercise}[2]
Prove that the sum of an even integer and an odd integer is odd.
\end{exercise}

\begin{exercise}[2]
Prove that if \(n\) is odd, then \(n^2\) is odd.
\end{exercise}

\begin{exercise}[2]
Use contraposition to prove:

If \(n^2\) is even, then \(n\) is even.
\end{exercise}

\begin{exercise}[2]
Use proof by contradiction to show that there is no greatest integer.
\end{exercise}

\begin{exercise}[2]
Prove by contradiction that \(\sqrt3\) is irrational.
\end{exercise}

\begin{exercise}[1]
Give a counterexample to the statement:

\begin{center}
For every real number \(x\),
\[ \sqrt{x^2}=x. \]
\end{center}
\end{exercise}

\begin{exercise}[1]
Give a counterexample to the statement:

\begin{center}
If \(ab=0\), then \(a=0\) and \(b=0\).
\end{center}
\end{exercise}

\begin{exercise}[2]
Prove that
\[ A\cap B\subseteq A. \]
\end{exercise}

\begin{exercise}[2]
Prove that
\[ A\subseteq A\cup B. \]
\end{exercise}

\begin{exercise}[2]
Prove that
\[ A\setminus B\subseteq A. \]
\end{exercise}

\begin{exercise}[2]
Prove
\[ A\cup B=B\cup A. \]
\end{exercise}

\begin{exercise}[3]
Prove
\[ A\cap(B\cup C)=(A\cap B)\cup(A\cap C). \]
\end{exercise}

\begin{exercise}[2]
Prove that the function
\[ f:\R\to\R,\qquad f(x)=7x+2 \]
is injective.
\end{exercise}

\begin{exercise}[2]
Prove that the same function is surjective.
\end{exercise}

\begin{exercise}[2]
Conclude that
\[ f(x)=7x+2 \]
is bijective.
\end{exercise}

\begin{exercise}[2]
Prove by induction that
\[ 1+3+5+\cdots+(2n-1)=n^2 \]
for every \(n\in\N\).
\end{exercise}

\begin{exercise}[2]
Prove by induction that
\[ 2^n\geq n+1 \]
for every \(n\in\N\).
\end{exercise}

\begin{exercise}[3]
Prove by induction that
\[ \sum_{k=1}^{n}k^2=\frac{n(n+1)(2n+1)}6. \]
\end{exercise}

\begin{exercise}[3]
Use strong induction to prove that every integer \(n\geq2\) is a product of
prime numbers.
\end{exercise}

\begin{exercise}[2]
Prove that for every integer \(n\), \(n(n+1)\) is even.
\end{exercise}

\begin{exercise}[3]
Prove that an integer \(n\) is odd if and only if \(n^2\) is odd.
\end{exercise}

\begin{exercise}[3]
Suppose \(f:A\to B\) and \(g:B\to C\) are injective. Prove that
\[ g\circ f \]
is injective.
\end{exercise}

\begin{exercise}[3]
Suppose \(f:A\to B\) and \(g:B\to C\) are surjective. Prove that
\[ g\circ f \]
is surjective.
\end{exercise}

%%%%%%%%%%%%%%%%%%%%%%%%%%%%%%%%%%%%%%%%%%%%%%%%%%%%%%%%%%%%
\subsection{Section Connections}
\label{subsec:proof-connections}
%%%%%%%%%%%%%%%%%%%%%%%%%%%%%%%%%%%%%%%%%%%%%%%%%%%%%%%%%%%%

\chapterconnection{
Proof methods are used throughout the remainder of mathematics. Direct proof,
contraposition, contradiction, and induction appear repeatedly in algebra,
number theory, analysis, topology, probability, and discrete mathematics.
Set proofs and function proofs provide patterns that will recur throughout
higher mathematics, while mathematical induction becomes especially important
for sequences, combinatorics, recurrence relations, algorithms, and number
theory.
}

%%%%%%%%%%%%%%%%%%%%%%%%%%%%%%%%%%%%%%%%%%%%%%%%%%%%%%%%%%%%
% SECTION SOURCE TRACKING
%%%%%%%%%%%%%%%%%%%%%%%%%%%%%%%%%%%%%%%%%%%%%%%%%%%%%%%%%%%%

%------------------------------------------------------------
% 0.8 Methods of Proof
%------------------------------------------------------------
%
% Sources:
%   - Richard Hammack, Book of Proof, Third Edition.
%     Key: hammack2018bookofproof
%
% Used for:
%   - direct proof structure;
%   - proof by contraposition and contradiction;
%   - existence and uniqueness arguments;
%   - counterexamples;
%   - biconditional proofs;
%   - set and function proof techniques;
%   - mathematical induction and strong induction;
%   - proof-writing conventions.
%
% Notes:
%   - Examples and exercises are independently developed.
%   - Formal proof theory is deferred to Chapter 12.
%   - Induction is revisited in discrete mathematics and number theory.
%
\nocite{hammack2018bookofproof}